% In this file you should put the actual content of the blueprint.
% It will be used both by the web and the print version.
% It should *not* include the \begin{document}
%
% If you want to split the blueprint content into several files then
% the current file can be a simple sequence of \input. Otherwise It
% can start with a \section or \chapter for instance.

\chapter{Nets}

\begin{definition}[Net and Marked Net]\label{def:net}
  \lean{net}
  \leanok
  \uses{def:net}
  A \emph{net} $N$ is a tuple $N = (\alpha, \beta, \preS{\_}, \postS{\_})$,
  where $\alpha$ is a nonempty set of places, $\beta$ is the set of
  transitions, and
  $\preS{\_}, \postS{\_}: \beta\rightarrow \nat^{\alpha}$ assign source and
  target to each transition. We consider only nets in which every transition
  both consumes and produces, i.e., $\preS{\tr{t}} \neq\emptyset$ and
  $\postS{\tr{t}} \neq\emptyset$ for all $\tr{t}\in \beta$.
  %
  A \emph{marking} $m$ of a net $N$ is a multiset of places, i.e.,
  $m\in\nat^{\alpha}$.
  %
  A \emph{marked net} is a pair $(N, m_0)$ where $N$ is a net and $m_0$ is
  a marking of $N$.
\end{definition}

The pre- and postsets of a place $p$ are $\preS{p} = \{t \mid p \in \postS{t}\}$ and $\postS{p} = \{t \mid p \in \preS{t}\}$.

\subsection{Operational semantics}

\begin{definition}[Enabled transition]\label{def:enabled}
  Let $N$ be a net. A transition $\tr{t}$ is \emph{enabled} at a marking $m$,
  denoted $m\llbracket \tr{t}\rrbracket_{\llBra N \rrBra}$, if $\preS{\tr{t}}\le m$.
\end{definition}

A transition is \emph{enabled} if the  marking contains its preset, i.e, as many tokens as those required by the transition's preset.
